\chapter{Complexity — DNF}
%%% generated by the blueprint pipeline from TCSlib.Complexity.Formulas.DNF
%%%%%%%%%%%%%%%%%%%%%%%%%%%%%%%%%%%%%%%%%%%%%%%%%%%%%%%%%%%%%%%%%%%%%%%%%%%%%%%
\section{Overview}

In [AB09, Section~2.6.1, Example~2.21] the Cook--Levin CNF formula is negated and one asks
whether the negation is a tautology; the negation of a CNF formula is a DNF formula, an OR
of ANDs of literals. This module supplies that dual layer: DNF formulas over variables
indexed by natural numbers, their evaluation and the tautology predicate, a binary
encoding shared with CNF formulas, the size-preserving De Morgan duals between CNF and DNF
formulas (which negate every literal in place), and the De Morgan bridge stating that the
dual of a CNF formula is a tautology exactly when the CNF formula is unsatisfiable.

Throughout, a literal is a pair $(v,b)$ of a variable index $v \in \mathbb{N}$ and a
polarity bit $b \in \{0,1\}$, satisfied by an assignment $a : \mathbb{N} \to \{0,1\}$
exactly when $a(v) = b$. A CNF formula is a finite list of clauses, each a finite list of
literals, true under $a$ when every clause contains a satisfied literal; a DNF formula is a
finite list of terms, each a finite list of literals, true under $a$ when some term has all
its literals satisfied.

%%%%%%%%%%%%%%%%%%%%%%%%%%%%%%%%%%%%%%%%%%%%%%%%%%%%%%%%%%%%%%%%%%%%%%%%%%%%%%%
\section{Declarations}

\begin{definition}[DNF formula]
\lean{Std.Sat.DNF}
\label{Std.Sat.DNF}
\leanok
A DNF formula over variables indexed by $\mathbb{N}$ is a finite list of terms
$(T_1, \dots, T_k)$, where each term $T_i$ is a finite list of literals $(v, b)$ with
variable index $v \in \mathbb{N}$ and polarity bit $b \in \{0,1\}$. Each term is read as
the conjunction of its literals and the formula as the disjunction of its terms; the
syntax is the same list-of-literal-lists shape as a CNF formula over $\mathbb{N}$, but it
carries this dual reading.
\end{definition}

\begin{definition}[Evaluation of a DNF formula]
\lean{Std.Sat.DNF.eval}
\label{Std.Sat.DNF.eval}
\leanok
\uses{Std.Sat.DNF}
Let $\psi = (T_1, \dots, T_k)$ be a DNF formula over variables indexed by $\mathbb{N}$,
each term $T_i$ a finite list of literals $(v,b)$ with $v \in \mathbb{N}$ and
$b \in \{0,1\}$, and let $a : \mathbb{N} \to \{0,1\}$ be an assignment. The value
$\psi(a) \in \{0,1\}$ is true exactly when some term $T_i$ has all of its literals
satisfied by $a$, a literal $(v,b)$ being satisfied when $a(v) = b$:
\[
  \psi(a) = \bigvee_{i=1}^{k} \; \bigwedge_{(v,b) \in T_i} \bigl[a(v) = b\bigr].
\]
In particular the formula with no terms evaluates to false under every assignment, and an
empty term is satisfied by every assignment [AB09, Section~2.6.1].
\end{definition}

\begin{definition}[Tautological DNF formula]
\lean{Std.Sat.DNF.Tautology}
\label{Std.Sat.DNF.Tautology}
\leanok
\uses{Std.Sat.DNF, Std.Sat.DNF.eval}
Let $\psi = (T_1, \dots, T_k)$ be a DNF formula over variables indexed by $\mathbb{N}$,
each term a finite list of literals $(v,b)$ with $v \in \mathbb{N}$ and $b \in \{0,1\}$.
The formula $\psi$ is a \emph{tautology} if for every assignment
$a : \mathbb{N} \to \{0,1\}$ some term $T_i$ has all of its literals satisfied by $a$,
where $(v,b)$ is satisfied when $a(v) = b$; that is, $\psi$ evaluates to true under every
assignment [AB09, Section~2.6.1].
\end{definition}

\begin{definition}[Decoding bit strings into DNF formulas]
\lean{Std.Sat.DNF.decode}
\label{Std.Sat.DNF.decode}
\leanok
\uses{Std.Sat.CNF.decode, Std.Sat.DNF}
For a bit string $x$, its DNF decoding is the DNF formula whose list of terms is the list
of clauses of the CNF decoding $\mathrm{dec}(x)$. Here $\mathrm{dec}(x)$ is the list of
literal lists obtained by parsing $x$, with recursion fuel $|x|$, according to the grammar
$\langle\text{formula}\rangle ::= 0 \mid 1\,\langle\text{clause}\rangle\,\langle\text{formula}\rangle$,
$\langle\text{clause}\rangle ::= 0 \mid \langle\text{literal}\rangle\,\langle\text{clause}\rangle$,
$\langle\text{literal}\rangle ::= 1^{k+1}\,0\,b$ (the literal with variable index
$k \in \mathbb{N}$ and polarity $b \in \{0,1\}$) when this parse consumes all of $x$, and
the empty list otherwise; so a malformed string decodes to the DNF formula with no terms.
\end{definition}

\begin{definition}[Binary encoding of a DNF formula]
\lean{Std.Sat.DNF.serialize}
\label{Std.Sat.DNF.serialize}
\leanok
\uses{Std.Sat.CNF.serialize, Std.Sat.DNF}
Let $\psi = (T_1, \dots, T_k)$ be a DNF formula over variables indexed by $\mathbb{N}$,
each term a finite list of literals $(v,b)$ with $v \in \mathbb{N}$ and $b \in \{0,1\}$.
Its binary encoding is the binary encoding of the same list of literal lists read as a CNF
formula,
\[
  \mathrm{enc}(\psi) = 1\,\mathrm{enc}(T_1)\;1\,\mathrm{enc}(T_2) \cdots 1\,\mathrm{enc}(T_k)\;0,
\]
where a term $(\ell_1, \dots, \ell_m)$ is encoded as
$\mathrm{enc}(\ell_1)\cdots\mathrm{enc}(\ell_m)\,0$ and a literal $(v,b)$ as $1^{v+1}0b$.
\end{definition}

\begin{theorem}[Decoding inverts the DNF encoding]
\lean{Std.Sat.DNF.decode_serialize}
\label{Std.Sat.DNF.decode_serialize}
\leanok
\uses{Std.Sat.CNF.decode_serialize, Std.Sat.DNF, Std.Sat.DNF.decode, Std.Sat.DNF.serialize}
\difficulty{1}
For every DNF formula $\psi$ over variables indexed by $\mathbb{N}$, decoding its binary
encoding as a DNF formula returns $\psi$. Here a DNF formula $(T_1,\dots,T_k)$, each term a
list of literals $(v,b)$ with $v \in \mathbb{N}$ and $b \in \{0,1\}$, is encoded as
$1\,\mathrm{enc}(T_1)\cdots1\,\mathrm{enc}(T_k)\,0$, a term as the concatenation of its
literal encodings $1^{v+1}0b$ followed by $0$; and a bit string $x$ is decoded as the DNF
formula whose terms are the literal lists obtained by parsing $x$, with recursion fuel
$|x|$, according to the grammar
$\langle\text{formula}\rangle ::= 0 \mid 1\,\langle\text{clause}\rangle\,\langle\text{formula}\rangle$,
$\langle\text{clause}\rangle ::= 0 \mid \langle\text{literal}\rangle\,\langle\text{clause}\rangle$,
$\langle\text{literal}\rangle ::= 1^{k+1}\,0\,b$ (the literal $(k,b)$) when this parse
consumes all of $x$, and as the formula with no terms otherwise.
\end{theorem}

\begin{theorem}[Decoding a CNF encoding as a DNF formula]
\lean{Std.Sat.DNF.decode_cnf_serialize}
\label{Std.Sat.DNF.decode_cnf_serialize}
\leanok
\uses{Std.Sat.CNF.decode_serialize, Std.Sat.CNF.serialize, Std.Sat.DNF.decode}
\difficulty{1}
Let $\varphi = (C_1, \dots, C_k)$ be a CNF formula over variables indexed by $\mathbb{N}$,
each clause a list of literals $(v,b)$ with $v \in \mathbb{N}$ and $b \in \{0,1\}$, with
binary encoding $\mathrm{enc}(\varphi) = 1\,\mathrm{enc}(C_1)\cdots1\,\mathrm{enc}(C_k)\,0$,
where a clause is encoded as the concatenation of its literal encodings $1^{v+1}0b$
followed by $0$. Then decoding $\mathrm{enc}(\varphi)$ as a DNF formula yields the DNF
formula whose terms are exactly $C_1, \dots, C_k$. Here a bit string $x$ is decoded as the
DNF formula whose terms are the literal lists obtained by parsing $x$, with recursion fuel
$|x|$, according to the grammar
$\langle\text{formula}\rangle ::= 0 \mid 1\,\langle\text{clause}\rangle\,\langle\text{formula}\rangle$,
$\langle\text{clause}\rangle ::= 0 \mid \langle\text{literal}\rangle\,\langle\text{clause}\rangle$,
$\langle\text{literal}\rangle ::= 1^{k+1}\,0\,b$ when this parse consumes all of $x$, and as
the formula with no terms otherwise.
\end{theorem}

\begin{definition}[De Morgan dual of a DNF formula]
\lean{Std.Sat.DNF.dual}
\label{Std.Sat.DNF.dual}
\leanok
\uses{Std.Sat.DNF}
Let $\psi = (T_1, \dots, T_k)$ be a DNF formula over variables indexed by $\mathbb{N}$,
each term a finite list of literals $(v,b)$ with $v \in \mathbb{N}$ and $b \in \{0,1\}$.
Its De Morgan dual is the CNF formula $(T_1', \dots, T_k')$, where the clause $T_i'$ is
obtained from $T_i$ by replacing every literal $(v,b)$ with $(v, \neg b)$, keeping the
order. Thus if $\psi = \bigvee_i \bigwedge_j \ell_{ij}$, its dual is the CNF formula
$\bigwedge_i \bigvee_j \neg\ell_{ij}$, of the same size.
\end{definition}

\begin{definition}[De Morgan dual of a CNF formula]
\lean{Std.Sat.CNF.dual}
\label{Std.Sat.CNF.dual}
\leanok
\uses{Std.Sat.DNF}
Let $\varphi = (C_1, \dots, C_k)$ be a CNF formula over variables indexed by $\mathbb{N}$,
each clause a finite list of literals $(v,b)$ with $v \in \mathbb{N}$ and
$b \in \{0,1\}$. Its De Morgan dual is the DNF formula $(C_1', \dots, C_k')$, where the
term $C_i'$ is obtained from $C_i$ by replacing every literal $(v,b)$ with $(v, \neg b)$,
keeping the order. Thus if $\varphi = \bigwedge_i \bigvee_j \ell_{ij}$, its dual is the DNF
formula $\bigvee_i \bigwedge_j \neg\ell_{ij}$; no distributive expansion takes place, and
the dual has the same size as $\varphi$.
\end{definition}

\begin{theorem}[Pointwise De Morgan law]
\lean{Std.Sat.CNF.eval_dual}
\label{Std.Sat.CNF.eval_dual}
\leanok
\uses{Std.Sat.CNF.dual, Std.Sat.DNF.eval}
\difficulty{4}
Let $\varphi = (C_1, \dots, C_k)$ be a CNF formula over variables indexed by $\mathbb{N}$,
each clause a finite list of literals $(v,b)$ with $v \in \mathbb{N}$ and
$b \in \{0,1\}$, and let $\varphi^{\ast}$ be the DNF formula whose terms are obtained from
$C_1, \dots, C_k$ by replacing every literal $(v,b)$ with $(v, \neg b)$. Then for every
assignment $a : \mathbb{N} \to \{0,1\}$,
\[
  \varphi^{\ast}(a) = \neg\, \varphi(a),
\]
where $\varphi(a)$ is true when every clause of $\varphi$ contains a literal $(v,b)$ with
$a(v) = b$, and $\varphi^{\ast}(a)$ is true when some term of $\varphi^{\ast}$ has all of
its literals $(v,b)$ satisfying $a(v) = b$.
\end{theorem}

\begin{theorem}[Dualizing a CNF formula twice]
\lean{Std.Sat.CNF.dual_dual}
\label{Std.Sat.CNF.dual_dual}
\leanok
\uses{Std.Sat.CNF.dual, Std.Sat.DNF.dual}
\difficulty{1}
Let $\varphi$ be a CNF formula over variables indexed by $\mathbb{N}$, a finite list of
clauses each a finite list of literals $(v,b)$ with $v \in \mathbb{N}$ and
$b \in \{0,1\}$. Forming its De Morgan dual DNF formula, by replacing every literal
$(v,b)$ with $(v,\neg b)$ and reading the clauses as terms, and then forming the De Morgan
dual CNF formula of that DNF formula, by again replacing every literal $(v,b)$ with
$(v,\neg b)$ and reading the terms as clauses, returns $\varphi$.
\end{theorem}

\begin{theorem}[Dualizing a DNF formula twice]
\lean{Std.Sat.DNF.dual_dual}
\label{Std.Sat.DNF.dual_dual}
\leanok
\uses{Std.Sat.CNF.dual, Std.Sat.DNF, Std.Sat.DNF.dual}
\difficulty{2}
Let $\psi$ be a DNF formula over variables indexed by $\mathbb{N}$, a finite list of
terms each a finite list of literals $(v,b)$ with $v \in \mathbb{N}$ and
$b \in \{0,1\}$. Forming its De Morgan dual CNF formula, by replacing every literal
$(v,b)$ with $(v,\neg b)$ and reading the terms as clauses, and then forming the De Morgan
dual DNF formula of that CNF formula, by again replacing every literal $(v,b)$ with
$(v,\neg b)$ and reading the clauses as terms, returns $\psi$.
\end{theorem}

\begin{theorem}[Dual tautology iff unsatisfiable]
\lean{Std.Sat.CNF.tautology_dual_iff}
\label{Std.Sat.CNF.tautology_dual_iff}
\leanok
\uses{Std.Sat.CNF.Satisfiable, Std.Sat.CNF.dual, Std.Sat.CNF.eval_dual, Std.Sat.DNF.Tautology}
\difficulty{1}
Let $\varphi = (C_1, \dots, C_k)$ be a CNF formula over variables indexed by $\mathbb{N}$,
each clause a finite list of literals $(v,b)$ with $v \in \mathbb{N}$ and
$b \in \{0,1\}$, and let $\varphi^{\ast}$ be the DNF formula whose terms are obtained from
$C_1, \dots, C_k$ by replacing every literal $(v,b)$ with $(v, \neg b)$. Then
$\varphi^{\ast}$ is a tautology, i.e.\ for every assignment $a : \mathbb{N} \to \{0,1\}$
some term of $\varphi^{\ast}$ has all of its literals $(v,b)$ satisfying $a(v) = b$, if and
only if $\varphi$ is unsatisfiable, i.e.\ no assignment makes every clause of $\varphi$
contain a literal $(v,b)$ with $a(v) = b$ [AB09, Section~2.6.1, Example~2.21].
\end{theorem}
